% !TeX root = ../main.tex
\subsection{Влияние амплитудно-фазового разбаланса (Пункт 4)}

В данном пункте исследуется влияние неидеальности формирования квадратурных составляющих гетеродина на спектр демодулированного сигнала. Моделирование проводится при фиксированной разрядности $n_g = 13$ бит (\texttt{nco\_type = 'single'}).

\begin{figure}[H]
    \centering
    \begin{minipage}[t]{0.48\textwidth}
        \vspace{0pt}
        \centering
        \includegraphics[width=\linewidth]{4_spec_a0_p0.png}
        \caption{Идеальные квадратуры}
        \label{fig:4_ideal}
    \end{minipage}
    \hfill
    \begin{minipage}[t]{0.48\textwidth}
        \vspace{0pt}
        \centering
        \includegraphics[width=\linewidth]{4_spec_a10_p0.png}
        \caption{Амплитудный разбаланс (1 дБ)}
        \label{fig:4_amv}
    \end{minipage}
\end{figure}

На рисунке \ref{fig:4_ideal} показан спектр при идеальных квадратурах (амплитудный разбаланс 0 дБ, фазовый $0^\circ$), где зеркальная составляющая полностью подавлена. На рисунке \ref{fig:4_amv} введено амплитудное искажение в 1 дБ, что приводит к нарушению симметрии ортогональных каналов и появлению побочного спектрального пика.

\begin{figure}[H]
    \centering
    \begin{minipage}[t]{0.48\textwidth}
        \vspace{0pt}
        \centering
        \includegraphics[width=\linewidth]{4_spec_a0_p5.png}
        \caption{Фазовый разбаланс ($5^\circ$)}
        \label{fig:4_pmv}
    \end{minipage}
    \hfill
    \begin{minipage}[t]{0.48\textwidth}
        \vspace{0pt}
        \centering
        \includegraphics[width=\linewidth]{4_spec_a15_p5.png}
        \caption{Разбаланс (1{,}5 дБ, $5^\circ$)}
        \label{fig:4_both}
    \end{minipage}
\end{figure}

Рисунок \ref{fig:4_pmv} иллюстрирует влияние исключительно фазовой ошибки в $5^\circ$, что также генерирует зеркальную помеху. В случае комплексного воздействия (рисунок \ref{fig:4_both}), включающего амплитудное рассогласование 1{,}5 дБ и фазовое $5^\circ$, мощность мешающей компоненты $P_r$ достигает максимального значения в рамках проведенных тестов, что согласуется с теоретической математической моделью демодулятора.
